% ==========================================================================
% GENERATED by scripts/split_paper.py. Do not edit this file.
%
% Source: imported_papers/probe-norm/content/body.tex, content/supplement.tex
% Part 9: The Published Models
%
% Edits belong upstream, in the published paper. To change what the
% chapter includes, change imports/bounded.toml and re-run
% the splitter.
% ==========================================================================
\section{The Published Models}
\label{app:published}
For every model we use the released checkpoints or the authors' code with their published settings, and fit our probe on encodings of held-out episodes (Table~\ref{tab:published_setup}). Table~\ref{tab:published_full} gives every diagnostic of the published models. We note only where we depart from the authors' protocol. Tables~\ref{tab:published} and~\ref{tab:published_more} collect each model's diagnostics at the horizon of its smallest factor, and Table~\ref{tab:fidelity} compares reproduced performance with published metrics.

\paragraph{Radius check.} Adding norm to direction changes held-out $R^2$ by at most $0.011$ for every published model, so re-projection loses no state in any of them.

\paragraph{Radius estimate.} Without an analytic radius, the reference radius is the norm of the nearest reference encoding in angle. To check that this choice does not drive the results, we recompute the factor at $T^*$ with half and a quarter of the reference bank, with the mean norm of the ten nearest encodings, and with the true target norm, which requires the future observation and serves only as a control. The estimates usually agree within a few percent (Table~\ref{tab:support}). Only the true target norm on TCRL's four late tasks gives a larger reduction. On predictions, the estimator's error can be measured in the trained modules against the analytic layer-normalization root. Nearest-encoding factors relative to analytic-root factors span $0.96$--$1.03$ at $T=5$ and $0.91$--$1.08$ at $T=20$. Large norm effects therefore do not depend on the estimator, while factors within a few percent of one do.

\input{tables/table_support.tex}

\input{tables/table_published_setup.tex}

\input{tables/table_published_full.tex}

\input{tables/table_fidelity.tex}

\subsection{LeWorldModel}
\label{app:lewm}
We use the released Push-T and Two-Room checkpoints with the authors' code. LeWorldModel's normalized output head bounds the norm without fixing a single radius. The Push-T checkpoint omits action normalization, so we select the action scale by one-step latent error on $40$ separate episodes, after which recorded actions outperform shuffled actions and persistence. Rollouts start from three encoded frames, and single-frame probes recover positions but not velocity. The Two-Room checkpoint is the Sub-JEPA baseline reported to show rollout drift~\citep{zhao2026subjepa}, and we evaluate it on expert-policy episodes that repeatedly cross the door.

On Push-T, norms remain stable through $T=500$, and position predictions beat both baselines through $T=20$. The architectural bound permits $2.1$ times the 99th-percentile encoding norm, so the observed stability is tighter than the bound. On Two-Room, the position probe has held-out $R^2=$ $0.9999$, which makes relative errors sensitive even to small norm deviations. Re-projection changes the one-step relative error from $0.016$ to $0.012$. From $T=20$, re-projection removes at most $11$\%, so norm mismatch explains little of the later error.

\subsection{TD-MPC2}
\label{app:tdmpc2}
We use three released checkpoints per task and verify our implementation against the official encoder and dynamics outputs. Short pendulum and acrobot evaluations use random actions. Cartpole and long-horizon evaluations use the checkpoint policy with Gaussian noise, because random actions visit unfamiliar cartpole states.

Pendulum and acrobot probes recover every state coordinate with held-out $R^2\ge$ $0.997$. Median norms remain within a few percent of the estimated reference radius through $T=400$. Relative error is $0.08$ at $T=50$ and exceeds one from $T=$ $150$. Re-projection to the exact SimNorm radius leaves all predictions unchanged, giving factor one (Section~\ref{app:radius}). The nearest-encoding estimator instead gives median factors of $0.94$--$1.09$ and individual values of $0.85$--$1.28$. Since the exact radius leaves every prediction unchanged, this variation is the estimator's error on predictions. TD-MPC2's long-horizon errors persist without norm mismatch.

\subsection{The Koopman autoencoder}
\label{app:kae}
We train ten plain and ten consistent models with the authors' code and settings~\citep{Azencot2020} and evaluate their own decoder and an auxiliary probe for pendulum angle and angular velocity. Rollouts start from held-out frames and run for $1000$ steps, $125$ times the training horizon.

One plain and three consistent runs have median norm ratios above ten at $T=1000$. We exclude them from retained-run summaries and separately count divergence as a loss in all-run comparisons (Table~\ref{tab:kae_seeds}). Among retained runs, target alignment remains high (median cosine at least $0.992$), while norms grow up to $3.2$-fold. Median re-projection factors are $0.14$ for the decoder and $0.021$ for the state probe. Using true target norms gives decoder factor $0.15$, close to the bank estimate. At $T=1000$, random directions with the predictions' norms give a state-probe factor of $0.83$, against $0.021$ for the predictions, and their error after re-projection exceeds the reference mean's, so the Koopman reduction reflects accurate directions.

Across retained runs, the rank correlation between decoder error and norm growth falls from $0.97$ to $0.14$ after re-projection. The consistent model wins $0.70$ of seed pairs before re-projection (interval $0.41$--$0.94$) and $0.38$ afterward ($0.11$--$0.68$). Under the authors' relative-error metric, its median error ratio to the plain model changes from $0.91$ to $1.80$, consistent with the change in decoder MSE. Excluding runs by norm growth conditions on the quantity being corrected, so Table~\ref{tab:kae_seeds} also reports comparisons over all runs and lists every run.

\input{tables/table_kae_seeds.tex}

\subsection{TCRL and TD-MPC}
\label{app:problem}
We train TCRL and TD-MPC with the authors' code and settings~\citep{zhao2023simplified,hansen2022temporal}, three seeds each, on pendulum, cartpole, and acrobot, and TCRL also on six further DeepMind Control tasks. Evaluation uses each model's policy with Gaussian action noise unless stated otherwise. Besides auxiliary state probes fitted on encodings, we read the reward heads the models train on predictions.

At $T=5$, TCRL's pendulum predictions have accurate directions under random actions (median target cosine $0.999$). Re-projection removes $97$\% of error under its policy and $89$\% under random actions. Under the policy, start-encoding probes achieve relative errors of $8.9\times10^{-4}$, $2.4\times10^{-3}$, and $3.8\times10^{-2}$ on pendulum, cartpole, and acrobot, and re-projected predictions give $2.3\times10^{-3}$, $1.8\times10^{-4}$, and $2.2\times10^{-3}$. Predictions improve on the start on cartpole and acrobot, while on pendulum the start-encoding probe is better by a factor of $2.6$. Probe fits disagree $22$ times more on re-projected pendulum predictions than on corresponding target encodings, but the absolute difference is only $2.4\times10^{-4}$ of the reference mean's error, compared with residual error $2.3\times10^{-3}$. TD-MPC factors remain near one through most tested horizons.

Other TCRL tasks show later reductions or none (Table~\ref{tab:published_full}), and random directions at the predictions' norms separate these late reductions from accurate prediction. TCRL predictions outperform random directions on all $27$ checkpoints at $T=5$. At each task's minimum-factor horizon, however, the pooled median factor is lower for random directions, and $18$ of $27$ prediction checkpoints still underperform the reference mean. These late reductions reflect the probe's sensitivity to scale.

We evaluate reward heads on raw predictions, re-projected predictions, and true future encodings against environment rewards. Median factors at $T=5$ are $1.03$ for TCRL and $1.08$ for TD-MPC, and $1.00$ and $1.00$ at $T=20$. Individual TCRL checkpoints range from $0.18$ to $2.09$ at $T=5$, with reductions on $9$ of $27$. At $T=20$, TCRL reward error is still $45$ times its error on target encodings.

\subsection{Semigroup-JEPA}
\label{app:sgjepa}
We reproduce the one-step baseline and the five-step Semigroup-JEPA with the authors' code on their square environment and keep their probe~\citep{Liu2026SemigroupJEPA}. The probe begins with layer normalization over a window of four latents and is therefore invariant to rescaling the whole window, apart from the small normalization stabilizer. We depart from this by estimating and restoring each predicted latent's radius separately before assembling the window, and we leave observed context latents unchanged. This can change relative scales within the window, so the probe's invariance to common scaling does not make it invariant to our intervention.

Individual rescaling changes probe output by a median $0.006$ target standard deviations. At $T=44$, the baseline's median norm ratios remain within $0.93$--$0.98$, and factors remain near one. Target directions are nearly orthogonal (median cosine $0.006$). Norm mismatch does not explain the baseline's error, and re-projection leaves the published comparison unchanged.

\subsection{KSL and PlayVirtual}
\label{app:ksl}
\label{app:playvirtual}
We train KSL~\citep{mcinroe2021learning} and PlayVirtual~\citep{yu2021playvirtual} from pixels with the authors' code, three seeds each on cartpole and cheetah-run. Both apply cosine consistency losses to projected latents, so we evaluate them both in raw latent space and in that projection space. After one cartpole step, raw norm ratios are $1.88$ and $2.51$, and re-projection removes $75$\% and $90$\% of probe error. On cheetah-run, held-out encoding $R^2$ falls to $0.26$, so we report these measurements in Table~\ref{tab:published_more} but do not interpret them as measures of rollout accuracy.

\paragraph{Projection space.} For KSL, PlayVirtual, and LightZero, we fit probes on projected target encodings and evaluate projected predictions (Table~\ref{tab:published_more}). KSL and PlayVirtual use moving-average target projection networks. LightZero shares one projection network between predictions and targets and fixes its batch-normalization statistics at evaluation. KSL and PlayVirtual have $50$-dimensional latents and projections, and LightZero projects $2304$ dimensions to $1024$. Each probe is fitted on $14880$ frames.

At $T=5$, all ten re-projected checkpoints outperform both baselines, although EfficientZero on Pong beats persistence by only $1.5$\%. Raw-space probes perform similarly for LightZero but remain far above the reference mean for KSL and PlayVirtual (Table~\ref{tab:published_more}). Re-projected random directions matched to prediction norms have relative error at least $1.9$ on every checkpoint. By $T=50$, only two checkpoints beat the reference mean, and median target cosine falls to $0.50$. Projected encodings also support less accurate state recovery, with median held-out $R^2$ of $0.69$ compared with $0.84$ in raw space, and $-0.04$ on MsPacman.

\input{tables/table_published_more.tex}

\subsection{EfficientZero and MuZero}
\label{app:lightzero}
We test seven released LightZero Atari checkpoints~\citep{niu2023lightzero} and retain six in Table~\ref{tab:published_more}: EfficientZero and Sampled EfficientZero on Pong and MsPacman, and MuZero on Breakout and Pong. LightZero trains MuZero with EfficientZero's consistency loss. Our episodes mix policy-network actions, without tree search, with random actions, and state targets are extracted from console memory~\citep{Anand2019AtariARI}. MuZero's Pong policy network always selects one action, so we collect episodes with EfficientZero's policy network and replay their actions through MuZero. MuZero on MsPacman is excluded because recorded actions give one-step latent error $4.93$, compared with $4.90$ for preceding actions and $2.46$ for persistence. MsPacman probes are also weak for EfficientZero and Sampled EfficientZero (median held-out $R^2$ of $0.31$ and $0.21$). EfficientZero's policy-prior return is substantially below its published search-based return (Table~\ref{tab:fidelity}).

Within training, factors lie in $0.97$--$1.07$. At $T=200$, norms reach $1.27$--$5.38$ times the reference radius, and re-projection removes $42$\%--$94$\% of error. All released configurations disable min--max scaling (\texttt{state\_norm}), and encoder and dynamics end in residual ReLU blocks with unbounded output norms.

\paragraph{Random-direction control.} At $T=200$, norm-matched random directions give factors $0.05$--$0.57$, within $0.16$ of prediction factors. On Pong, target encodings assigned predicted norms give factors $0.011$--$0.06$ and recover their original error after re-projection. On MsPacman, prediction norms are two to five times target norms, and factors for predictions, random directions, and rescaled targets coincide.
%%END_SUPPLEMENT
